\documentclass[12pt]{article}

\usepackage{graphicx}
\usepackage{hyperref}
\usepackage{geometry}
\usepackage{float}

\geometry{margin=1in}

\title{Chest X-Ray Pneumonia Detection Using Deep Learning}
\author{Maulik Gajera}
\date{\today}

\begin{document}

\maketitle

\begin{abstract}
Pneumonia is a serious lung infection that can be detected using chest X-ray imaging.
This project applies deep learning techniques to automatically classify chest X-ray
images as Normal or Pneumonia. Multiple convolutional neural network models are
implemented and compared to evaluate their effectiveness for medical image classification.
\end{abstract}

\section{Introduction}

Medical imaging plays an important role in diagnosing diseases. Chest X-rays are
commonly used to detect lung infections such as pneumonia. However, manual
interpretation by radiologists can take time and may vary depending on experience.

Deep learning techniques, particularly Convolutional Neural Networks (CNNs),
have shown strong performance in image classification tasks. This project
explores the use of CNN-based architectures to detect pneumonia from chest
X-ray images.

\section{Problem Statement}

The objective of this project is to build a deep learning model capable of
classifying chest X-ray images into two categories:

\begin{itemize}
\item Normal
\item Pneumonia
\end{itemize}

\section{Dataset}

The dataset used in this project is the Chest X-Ray Pneumonia dataset.

Dataset structure:

\begin{verbatim}
chest_xray/
    train/
        NORMAL/
        PNEUMONIA/
    test/
        NORMAL/
        PNEUMONIA/
    val/
        NORMAL/
        PNEUMONIA/
\end{verbatim}

The dataset contains labeled X-ray images used for training and evaluation.

\section{Data Preprocessing}

The following preprocessing steps were applied:

\begin{itemize}
\item Image resizing
\item Normalization
\item Data augmentation
\item Train / validation / test split
\end{itemize}

These steps improve model performance and generalization.

\section{Models Implemented}

Several deep learning models were trained and evaluated:

\subsection{Custom CNN}
A baseline convolutional neural network built using multiple convolution
and pooling layers.

\subsection{ResNet}
Residual networks help train deeper models and improve feature extraction.

\subsection{MobileNet}
A lightweight architecture optimized for efficiency and speed.

\subsection{EfficientNet}
A modern architecture balancing depth, width, and resolution.

\section{Training Process}

The models were trained using GPU acceleration with the following configuration:

\begin{itemize}
\item Optimizer: Adam
\item Loss Function: CrossEntropyLoss
\item Batch Size: 32
\item Epochs: 10 -- 20
\end{itemize}

\section{Evaluation Metrics}

Model performance was evaluated using:

\begin{itemize}
\item Accuracy
\item Precision
\item Recall
\item Confusion Matrix
\end{itemize}

\section{Results}

\begin{table}[H]
\centering
\begin{tabular}{|c|c|}
\hline
Model & Accuracy \\
\hline
CNN & -- \\
ResNet & -- \\
MobileNet & -- \\
EfficientNet & -- \\
\hline
\end{tabular}
\caption{Model Comparison}
\end{table}

\section{Visualization}

\begin{figure}[H]
\centering
\includegraphics[width=0.7\linewidth]{images/accuracy.png}
\caption{Training Accuracy}
\end{figure}

\begin{figure}[H]
\centering
\includegraphics[width=0.7\linewidth]{images/loss.png}
\caption{Training Loss}
\end{figure}

\section{Conclusion}

This project demonstrates that deep learning models can effectively detect
pneumonia from chest X-ray images. CNN-based architectures are capable of
learning meaningful patterns from medical imaging data and assisting in
automated diagnosis.

\section{Future Work}

Future improvements may include:

\begin{itemize}
\item Larger and more diverse datasets
\item Hyperparameter optimization
\item Model explainability using Grad-CAM
\item Deployment as a web application
\end{itemize}

\section{References}

\begin{itemize}
\item Kaggle Chest X-Ray Dataset
\item Deep Learning for Medical Image Analysis
\item Convolutional Neural Networks Research Papers
\end{itemize}

\end{document}